\section{A path-local theorem for ordinary flops}
\label{sec:ordinary-flop}

Let \(Y\dashrightarrow Y'\) be a projective ordinary flop covered by
\cite{LeeLinQuWang}, and let \(\ell,\ell'\) be its effective extremal fibre
classes.  The graph correspondence \(F\) identifies the big quantum products
and Poincar\'e pairings after analytic continuation
\begin{equation}\label{eq:flop-novikov-inversion}
 q'=q^{-1},\qquad F(\ell)=-\ell'.
\end{equation}
The continuation is analytic in the extremal variable and formal in the
other Novikov variables.

Fix a simply connected nonsingular domain \(U\) for the continued extremal
variable.  We use the hybrid coefficient ring
\begin{equation}\label{eq:flop-coefficient-ring}
 \mathcal O(U)\widehat\otimes\Lambda_{\perp}[[\tau]].
\end{equation}
where \(\Lambda_{\perp}\) is the transverse Novikov completion and \(\tau\)
denotes the bulk variables; the coefficient recursion below is taken after
extension to \(\C((z))\).  The \emph{graph gauge} is the horizontal gauge
obtained from the Lee--Lin--Qu--Wang comparison of big quantum connections on
\(U\), normalized by the classical graph correspondence.  Thus it acts on
flat columns, not only on quantum products.

Choose corresponding points \(p\in Y\) and \(p'\in Y'\) outside the
exceptional loci.  Let \(A\) be an ample divisor on the common contraction
and write \(D,D'\) for its pullbacks.  Then
\(D\cdot\ell=D'\cdot\ell'=0\), while \(D\) is positive on every effective
class away from the extremal ray.  We complete the transverse Novikov ring
by \(D\)-degree.

\begin{theorem}[Point-section preservation]
\label{thm:ordinary-flop-point-row}
On \(U\), the graph gauge satisfies
\begin{equation}\label{eq:flop-point-section}
 F(s_Y(\OO_p))=s_{Y'}(\OO_{p'}).
\end{equation}
Consequently it preserves the Gamma point row.
\end{theorem}

\begin{proof}
At transverse degree zero, every nonconstant stable map has pure extremal
class and lies in the exceptional locus.  It cannot meet the chosen point.
All positive-extremal descendant invariants with the point insertion
therefore vanish.  The point columns on both sides are their common
classical column for every nonsingular value of \(q\):
\(\Ghat_Y\operatorname{ch}(\OO_p)\) is a scalar multiple of the top-degree
class \([p]\), on which higher Gamma and Chern-class terms vanish.  Analytic
continuation of this identically classical extremal column remains
classical, and the graph correspondence sends \([p]\) to \([p']\).

Let
\[
 \delta=F(s_Y(\OO_p))-s_{Y'}(\OO_{p'}).
\]
It is jointly flat, meaning flat for both the \(z\)- and Novikov-direction
connections on the hybrid ring \eqref{eq:flop-coefficient-ring}, and has zero
transverse constant term.  Suppose
\(\delta_\beta\) is a nonzero coefficient of minimal positive \(D\)-degree.
Take coefficients modulo the extremal ray.  By minimality, every
positive-transverse convolution term involves a lower \(D\)-degree
coefficient and vanishes.  Horizontality for the Novikov derivation defined
by \(D\) therefore gives
\begin{equation}\label{eq:flop-transverse-tail}
 \bigl((D\cdot\beta)\id+z^{-1}(D\star_q-)\bigr)
 \delta_\beta=0.
\end{equation}
At transverse degree zero, quantum corrections to \(D\star_q\) have pure
extremal degree.  The divisor axiom multiplies them by
\(D\cdot\ell=0\).  Thus \(D\star_q=D\cup-\) in
\eqref{eq:flop-transverse-tail}.  Cup product by \(D\)
is nilpotent and \(D\cdot\beta>0\), so the displayed operator is invertible
over \(\C((z))\).  This contradiction kills the first possible coefficient.
Induction proves \eqref{eq:flop-point-section}.  Flatness of the pairing then
preserves the row \eqref{eq:gamma-point-row}.
\end{proof}

\begin{corollary}[Packetwise Boolean]
\label{cor:ordinary-flop-packet}
Let \(\cP\subset\mathcal S(Y)\) be a sectorially realized formal-monodromy
packet whose chosen realization is transported by the Lee--Lin--Qu--Wang
graph gauge to a sectorially realized packet
\(\cP'\subset\mathcal S(Y')\).  On the fixed continuation domain,
\begin{equation}\label{eq:flop-primary-boolean}
 b(Y,\cP)=b(Y',\cP').
\end{equation}
\end{corollary}

\begin{proof}
The graph gauge preserves the quantum connection, grading, first Chern class,
and Poincar\'e pairing.  It therefore transports the complete \(z=0\) formal
type.  Theorem~\ref{thm:ordinary-flop-point-row} identifies the covector on
the transported sectorial realization, and
\eqref{eq:flop-primary-boolean} follows.
\end{proof}

The theorem is narrower than a Gamma/Fourier--Mukai comparison for every
\(K\)-class.  Pure extremal curves need not miss the support of a general
class, and the support computation in the first paragraph of the proof would
fail.  Nor does the argument assert phase independence outside the chosen
continuation domain.
